% !TeX root = ../main.tex
% Placeholder reused by every check of this section. Drop the screenshot at
% figures/simValidationCaptureRun.png and all the figures of the section update at
% once; while the file is missing, the example-image shipped with the mwe package is
% typeset instead so the document keeps compiling.
\newcommand{\physfig}[2]{%
  \par\vspace{6pt}%
  \begin{center}%
    \IfFileExists{figures/#1.png}{%
      \includegraphics[width=0.85\linewidth]{#1}%
    }{%
      \includegraphics[width=0.85\linewidth]{example-image}%
    }\\[2pt]%
    \captionof{figure}{#2}\label{fig:#1}%
  \end{center}%
  \vspace{2pt}%
}

The simulator described in Section~\ref{sec:training_set} replaces heuristic signal
generation with an explicit powertrain model: an armature circuit, a torque balance
referred to the motor shaft, a lumped thermal node for the winding, a battery with an
equivalent-circuit terminal voltage and a set of transducers that inject noise,
quantisation and drift into every published channel. Introducing that model raises a
question that unit tests alone cannot answer: whether the telemetry that finally
reaches the historian still \emph{looks} like the physical system it claims to
represent.

\textbf{Nature of this test.} The checks collected in this section are
\emph{functional, informational tests} executed by \textbf{manual inspection} of the
telemetry rendered in the Grafana dashboards. They are not automated tests, they are
not executed by a CI/CD pipeline, they produce no pass/fail artefact and they are not
regression protected: every run is observed, interpreted and archived by hand. Their
purpose is evidential rather than verificational, to show a reader of the thesis that
the simulated signals reproduce the qualitative and quantitative behaviour expected
from the electromechanical model, so that the anomalies synthesised later on can be
argued to be physically plausible and not artefacts of arbitrary signal shaping. The
deterministic arithmetic of the model (torque equilibrium, back-EMF, energy balance,
sensor conditioning, fault signatures) is instead covered by the automated Rust test
suite of the simulator, which is executed on every build and is the appropriate
instrument for that class of assertions. This section therefore complements, and does
not replace, the automated tests.

\subsection{Test Environment and Procedure}

All observations were taken from the deployed stack: the simulated fleet publishes one
VDA~5050 \texttt{state} message per robot and second, which travels through the MQTT broker
through Telegraf and Kafka Connect into the QuestDB table \texttt{vda5050\_state}, where
the Grafana dashboards \textit{VDA5050 Powertrain Physics}, \textit{VDA5050 Injected
Faults}, \textit{VDA5050 Robot Dashboard}, \textit{VDA5050 Robot Velocity} and
\textit{VDA5050 Robot Vibration} read them. The state publication rate of one hertz is
the effective sampling rate of every observation reported below and sets the limit of
what can be seen.

The reference stimulus is a single long mission issued with the \textbf{Capture run
(long straight)} control-panel action, which builds a straight unloaded route towards
the positive x direction with a node every 60 m, sized from the configured cruise speed
so that the vehicle remains at cruise for roughly eleven minutes. A straight, unloaded
path is deliberate: it keeps the lateral and angular velocity at zero, so the only
events visible in the plots are the repeated acceleration ramps produced by the nodes
and the slower thermal and electrochemical ramps that accumulate behind them. Each
check below states what has to be observed and the value measured in the reference run;
the figure that follows every check is the corresponding panel of the dashboard.

Two conventions have to be kept in mind when reading the numbers, because without them
two of the published channels appear inconsistent with the model. First, the current
transducer is a hall-effect sensor clamped on the DC side of the drive, so the published
\texttt{motor\_current} is the armature current plus the standby draw of the drive
itself, denoted $I_{\mathrm{q}} = 0.35~\mathrm{A}$; the torque-producing current is
therefore $I_{\mathrm{m}} = \texttt{motor\_current} - I_{\mathrm{q}}$, and the torque
channel is consistent with the current channel only after that subtraction. Second, the
published \texttt{motor\_voltage} is the DC-link, that is, the battery terminal
voltage, and not the voltage that the converter applies to the armature, which is the
lower quantity appearing in Equation~\eqref{eq:motor_voltage}; combining the published
voltage with the published back-EMF and armature resistance would otherwise seem to
violate that equation by about four volts. Throughout this section, \emph{armature
current} refers to \texttt{motor\_current} and \emph{pack current} refers to
\texttt{battery\_current\_a}, which are two different channels.

\subsection{Start-Up and Cruise}

\begin{enumerate}
    \item \textbf{Current peak during acceleration, followed by a cruise plateau.}
    Issuing the order produces an armature current, \texttt{motor\_current}, of
    approximately $8.6~\mathrm{A}$ lasting two seconds, after which the signal settles
    on a plateau of about $3.3~\mathrm{A}$ that it does not leave while the vehicle
    travels at cruise. The pack current, \texttt{battery\_current\_a}, traces the
    same profile at a lower amplitude, the ratio between the two channels not being
    constant because it depends on the voltage the converter applies to the armature, as
    discussed in the battery subsection. The peak is the accelerating
    load of Equation~\eqref{eq:torque_balance}: the vehicle goes from $0$ to
    $1.0~\mathrm{m/s}$ at $0.5~\mathrm{m/s^2}$, and both the inertial and the rolling
    terms of the load torque are paid for with armature current. The plateau is the same
    equation with the inertial term cancelled. This two-second profile is not
    the electrical transient of the machine: the armature circuit has an inductance and a
    resistance whose quotient, $L_{\mathrm{a}}/R_{\mathrm{a}} \approx 40~\mathrm{ms}$,
    is the time the current needs to reach the value demanded by the drive. That interval
    is twenty-five times shorter than the one second between samples, so the transient is
    completed inside a single sample and cannot appear in the plot; its absence is a
    property of the sampling rate, not evidence that the armature inductance is missing
    from the model.

    \physfig{armature_current_and_bus_voltage_arranque_crucero}{Armature current at start-up: an approximately
    $8.6~\mathrm{A}$ peak during the two-second acceleration ramp followed by a
    $3.3~\mathrm{A}$ cruise plateau.}

    \item \textbf{Motor speed reaching the value imposed by the transmission.}
    The measured speed goes from $0$ to $100~\mathrm{rad/s}$ in the same two seconds and
    then stays flat. The value of the plateau is the relevant evidence, not its shape:
    it is $\omega_{\mathrm{m}} = i_{\mathrm{g}}\, v / r_{\mathrm{w}} = 10 \times
    1.0~\mathrm{m/s}\,/\,0.10~\mathrm{m} = 100~\mathrm{rad/s}$: the wheel turns at
    $10~\mathrm{rad/s}$ and the gearbox multiplies that speed by ten, which is the
    kinematic relation stated in Section~\ref{sec:training_set}. A plateau at any other
    value would mean that the wheel radius or the gearbox ratio published in the
    configuration is not the one the model uses.

    \physfig{motor_speed_back_EMF_arranque_crucero}{Motor angular velocity and back-EMF during start-up and cruise: the speed reaches
    $100~\mathrm{rad/s}$ in two seconds, and the back-EMF reaches $19.00~\mathrm{V}$ during cruise.}

    \item \textbf{Back-electromotive force proportional to speed.}
    As soon as the speed reaches its plateau, \texttt{back\_emf\_v} sticks to
    $19.00~\mathrm{V}$. The check consists of dividing the reported quantity by
    $K_{\mathrm{e}} \omega_{\mathrm{m}}$: the measured ratio is $1.0001$, that is,
        Equation~\eqref{eq:motor_back_emf} holds with the configured constant and is not
    being emulated from a different quantity. This observation is supported by the same
    panel used for the speed check, \Cref{fig:motor_speed_back_EMF_arranque_crucero}.

    \item
 \textbf{Torque consistent with current after removing the standby draw.}
    The reported torque has to satisfy $\tau_{\mathrm{m}} = K_{\mathrm{t}}
    (I_{\mathrm{m}} - I_{\mathrm{q}})$; at the measured cruise current of
    $3.33~\mathrm{A}$ this yields $0.57~\mathrm{N\,m}$, which is the value on the plot.
    Subtracting $I_{\mathrm{q}}$ is not a correction applied after the fact: the
    Hall-effect transducer of the model sees the standby consumption of the drive
    together with the armature current, so a torque computed from the raw current would
    be systematically biased. This is the discrete counterpart of
    Equation~\eqref{eq:electromagnetic_torque}.

    \physfig{motor_torque_crucero}{Electromagnetic torque against armature current. The slope
    matches $K_{\mathrm{t}}$ once the $0.35~\mathrm{A}$ standby draw of the drive is
    removed from the measured current.}

    \item \textbf{Shaft power equal to torque times speed.}
    The power channel reads about $57~\mathrm{W}$ at cruise, which is
    $0.57~\mathrm{N\,m} \times 100~\mathrm{rad/s}$. This closes the energy chain between
    the mechanical and the electrical channels and rules out that power is generated
    independently from torque and speed, as shown in
    \Cref{fig:motor_torque_crucero}.

    \item \textbf{Current is a band, not a line.}
    The cruise armature
    current, \texttt{motor\_current}, is a band of roughly $\pm 0.08~\mathrm{A}$ around
    its $3.3~\mathrm{A}$ plateau; the standard deviation measured over the reference run
    is $0.082~\mathrm{A}$ against the $0.08~\mathrm{A}$ configured for the current
    transducer, the small excess being the twelve-bit quantisation of an eighty-ampere
    range, one least significant bit being $19.5~\mathrm{mA}$. This is the evidence that the sensor
    conditioning layer (Gaussian noise, analogue-to-digital conversion and drift) is
    active on the published signal. Two opposite defects are ruled out at once: a band
    clearly wider than configured, which would indicate that fault ripple or a wrong
    noise scale is leaking into healthy data, and a perfectly straight line, which would
    mean the conditioning layer has been bypassed and the training set would be
    unrealistically clean, as observed in
    \Cref{fig:armature_current_and_bus_voltage_arranque_crucero}.

    \item \textbf{Vibration rises with travel, not with command.}
    The channel reports the root-mean-square acceleration of the chassis in units of
    $g$, the scalar an accelerometer with an internal RMS stage publishes for each
    measuring window; the waveform itself does not travel in the state message. In the
    model that scalar is driven by the quantities that excite the chassis: an idle floor
    of $0.02~\mathrm{g}$, a rolling term of $0.15~\mathrm{g}$ per $\mathrm{m/s}$ of
    speed actually made good over the ground, the fault-driven terms described in
    Section~\ref{sec:training_set} and the noise of the transducer, with a dispersion of
    $0.025~\mathrm{g}$. At the cruise speed of $1.0~\mathrm{m/s}$ the expected value is
    $0.02 + 0.15 \times 1.0 = 0.17~\mathrm{g}$, about $1.7~\mathrm{m/s^2}$, and the
    reading has to be interpreted as a distribution around that mean rather than as a
    level: single samples between $0.12$ and $0.22~\mathrm{g}$ are normal, and a value
    of $0.206~\mathrm{g}$ corresponds to $1.4$ standard deviations above the mean, which
    is expected in roughly one sample out of thirteen. The assertion is therefore made on
    the mean level and on the timing, not on individual samples. The level jumps from a
    $0.02~\mathrm{g}$ idle floor to a $0.17~\mathrm{g}$ mean with the same two-second
    delay as the speed, never before it, because the excitation of the model is the speed
    made good over the ground and not the velocity command: an AGV whose wheels are
    commanded but which does not advance must not show cruising vibration.

    \physfig{vibration_rise_arranque}{Chassis RMS vibration: $0.02~\mathrm{g}$ idle
    floor at standstill and a $0.17~\mathrm{g}$ mean, with the $0.025~\mathrm{g}$
    dispersion of the transducer, while travelling; the rise follows the speed.}
\end{enumerate}

\subsection{Thermal Behaviour}

\begin{enumerate}
    \item \textbf{Winding temperature rises with a curve towards a predictable
    equilibrium.}
    The winding starts at the ambient temperature of $22.0~^{\circ}\mathrm{C}$ and climbs
    to $35.1~^{\circ}\mathrm{C}$ along the seven hundred seconds of travel of the
    reference run. The quantity that has to be evaluated is neither the average slope,
    which depends on how much of the curve was recorded, nor the peak value, but the pair
    formed by the time constant and the asymptote. Fitting the response of a first-order
    node, $T(t) = T_{\infty} + (T_0 - T_{\infty})\,e^{-(t-t_0)/\tau_{\mathrm{th}}}$,
    to the samples stored in the historian yields $\tau_{\mathrm{th}} = 211~\mathrm{s}$
    and $T_{\infty} = 35.6~^{\circ}\mathrm{C}$, with a residual dispersion of
    $0.05~\mathrm{K}$, which is the noise of the thermistor and nothing else. Both fitted
    parameters correspond to an expression of the model: the time constant is
    $R_{\mathrm{th}} C_{\mathrm{th}} = 2.5 \times 80 = 200~\mathrm{s}$, and the
    asymptote is the ambient temperature plus the dissipated power times the thermal
    resistance, $22.0 + 5.4 \times 2.5 = 35.5~^{\circ}\mathrm{C}$, for the
    $0.45~\mathrm{W}$ of copper loss and $5.0~\mathrm{W}$ of iron loss of this operating
    point. The criterion that can be applied by hand on the panel is the sixty-three
    percent rule: the temperature has to have covered $0.63 \times 13.6 = 8.6~\mathrm{K}$
    of the $13.6~\mathrm{K}$ span after one time constant, which means crossing
    $30.6~^{\circ}\mathrm{C}$ two hundred seconds after the drive is loaded. A rise that
    were linear, or that reached its plateau within one or two samples, would invalidate
    the thermal node.

    \physfig{temperature_meseta}{Winding temperature during the capture run: exponential
    rise from the ambient $22.0~^{\circ}\mathrm{C}$ to $35.1~^{\circ}\mathrm{C}$, with
    a fitted time constant of $211~\mathrm{s}$ and a fitted asymptote of
    $35.6~^{\circ}\mathrm{C}$ against the $35.5~^{\circ}\mathrm{C}$ predicted by
    $T_{\mathrm{amb}} + P R_{\mathrm{th}}$.}

    \item \textbf{Cooling with the same exponent, not with the mirror curve.}
    Heating and cooling are not mirror images and must not be compared by eye: the heating
    branch approaches the equilibrium of the operating point, $35.5~^{\circ}\mathrm{C}$,
    whereas the cooling branch approaches the ambient $22.0~^{\circ}\mathrm{C}$, where
    no loss is being generated, so the two have different asymptotes and the piece
    of each one that appears on the panel depends on how long the vehicle ran. What has to
    be identical in both directions is the exponent. Once the order is finished, the
    temperature falls from $35.1$ to $22.15~^{\circ}\mathrm{C}$ in twenty-two minutes,
    and the same fit returns $\tau_{\mathrm{th}} = 208~\mathrm{s}$ against the
    $211~\mathrm{s}$ of the heating branch, with the same $0.05~\mathrm{K}$ residual: one
    single lumped node losing heat through a fixed resistance heats and cools with the same
    time constant in both directions, as it must. Without fitting, the equivalent criterion
    is that after one time constant the cooling has covered $63~\%$ of the distance to the
    ambient, near $26.9~^{\circ}\mathrm{C}$, and that the logarithm of that distance is a
    straight line of slope $-1/\tau_{\mathrm{th}}$ on both branches. A temperature
    collapsing to the ambient in one or two seconds would mean that the thermal state is
    being overwritten instead of integrated.

    \physfig{temperature_godown}{Cooling of the winding after the mission ends: from
    $35.1$ to $22.15~^{\circ}\mathrm{C}$ with a fitted time constant of
    $208~\mathrm{s}$, indistinguishable from the $211~\mathrm{s}$ of the heating branch,
    and with a residual equal to the noise of the thermistor.}

    \item \textbf{The battery node is deliberately sluggish.}
    The pack temperature barely moves over several minutes, by a few tenths of a degree.
    This is expected and must not be mistaken for a defect: the pack node has a heat
    capacity of $900~\mathrm{J/K}$ against $80~\mathrm{J/K}$ for the winding and is
    heated only by $I_{\mathrm{bat}}^2 R_{\mathrm{int}}$, which is small at cruise. Its
    slow drift is nevertheless relevant because the internal resistance of the battery
    depends on temperature.

    \physfig{simval-thermal-battery}{Winding and pack temperature on the same axis: the
    pack node responds one order of magnitude more slowly than the winding.}
\end{enumerate}

\subsection{Braking and Node Arrival}

\begin{enumerate}
    \item \textbf{Current falls to the standby value during braking, never below it.}
    As the vehicle approaches a node, the reported longitudinal velocity ramps down and
    the armature current drops from $3.3~\mathrm{A}$ to about $0.33~\mathrm{A}$, the
    standby draw of the drive, in a single sample. The sign is what has to be checked:
    the current must not become negative. The model does not implement regenerative
    braking, so deceleration dissipates energy mechanically and the drive simply stops
    demanding torque. A negative current on this plot would therefore be a defect, not a
    feature.

    \physfig{current_goes_down_break}{Velocity and armature current at node arrival: the
    current returns to the standby level and never crosses into negative values.}

    \item \textbf{Vibration and DC-link voltage recover in the same second.}
    The vibration returns to the $0.02~\mathrm{g}$ idle floor, mean value and same
    dispersion as above, in the same second in which the velocity reaches zero, because
    the excitation of the accelerometer model is the speed made good over the ground. The
    DC-link voltage also steps back up as soon as the load is removed, but this second
    recovery is not legible on the raw trace and should not be looked for there: the
    magnitude of the step is $120~\mathrm{mV}$, whereas the accelerated discharge moves
    the open-circuit voltage of the pack by $53~\mathrm{mV}$ for every minute of travel,
    in the same direction before and after the stop, so the load step is buried under a
    trend four times faster than itself.

    \physfig{vibration_reach_default_break}{Vibration and DC-link voltage at standstill
    recovery. The vibration collapses to its idle floor in one sample; the voltage
    recovery of $120~\mathrm{mV}$ is masked by the $53~\mathrm{mV}$ per minute decline of
    the open-circuit voltage.}
\end{enumerate}

\subsection{Battery}

\begin{enumerate}
    \item \textbf{State of charge decreases linearly.}
    The state of charge follows a straight ramp downwards of about two percent per
    minute; the reference run consumed $7.5~\%$ in three minutes. The slope is constant
    because the cruise current is constant and the charge is integrated by Coulomb
    counting, and because the discharge rate used in the demonstration configuration is
    accelerated by an intentional factor of twenty to make a mission-sized drain visible
    in a few minutes of observation. That factor multiplies elapsed time only, so it
    changes how fast the ramp runs but not its linearity, and it must be reported
    together with the slope whenever the value is quoted.

    \physfig{battery_linear_discharge}{State of charge during the capture run: linear discharge ramp,
    measured with the accelerated demo discharge factor enabled.}

    \item \textbf{Battery current is lower than armature current.}
    The pack delivers $2.90~\mathrm{A}$ while the armature carries $3.33~\mathrm{A}$.
    The ratio of about $0.8$ is the consequence of the converter: the armature is fed at
    a voltage lower than the DC-link, so the current drawn from the pack is reduced by
    approximately the same factor; the residual difference is the auxiliary consumption
    of the on-board electronics and the standby draw of the drive, $0.50~\mathrm{A}$ in
    the configuration used, which flows from the pack without passing through the
    armature. The ratio is therefore load dependent: it approaches one at high speed,
    when the converter applies nearly the full bus voltage to the armature, and it
    becomes much smaller during a start-up, when the armature voltage is low. A ratio close to one, or a battery current larger than the
    armature current while travelling, would break the energy balance of the model.

    \physfig{battery_current_lower_motor_current}{Battery current against armature current during
    cruise; the ratio reflects the converter voltage transformation plus the auxiliary
    consumption.}

    \item \textbf{Terminal voltage is consistent with the equivalent circuit.}
    Equation~\eqref{eq:battery_ecm} implies a drop of about $150~\mathrm{mV}$ from the
    open-circuit value at the cruise current of $2.96~\mathrm{A}$, and therefore a step of
    about $120~\mathrm{mV}$ between the loaded and the unloaded level. That step is not
    legible on the raw trace and no figure is offered for it: the accelerated discharge of
    the demonstration configuration drags the open-circuit voltage downwards at
    $53~\mathrm{mV}$ per minute of travel, four times faster than the step itself and
    undisturbed by the switching of the drive, so recovering the identity from the panels
    would require compensating the state of charge analytically. The arithmetic of the
    equivalent circuit, including the relaxation of the polarisation branch with the time
    constant $R_{\mathrm{p}} C_{\mathrm{p}} = 300~\mathrm{s}$, is asserted in the
    automated test suite of the simulator, which exercises the battery model directly, and
    is deliberately left out of the manual inspection reported here.
\end{enumerate}

\subsection{Interpretation and Limitations}

The values quoted above are consistency checks against the model, not metrological
measurements: every one of them is conditioned by the calibration constants published in
the simulator configuration, and they hold for the operating point described, a
one-hundred-and-eighty-kilogram vehicle travelling unloaded on level ground at
$1.0~\mathrm{m/s}$. Changing the payload, the grade or the speed moves the numerical
values while leaving the relations intact, which is precisely the property the checks
exercise.

The one hertz publication interval also bounds the scope of the section. Electromagnetic
transients of tens of milliseconds, the ripple harmonics injected by the stator
short-circuit fault, and the slow drift of the thermistor and the current transducer,
specified in fractions of a degree or a percent per hour, are outside what can be
observed here; they are asserted deterministically in the automated test suite of the
simulator instead. This section is consequently an informational, human-reviewed
characterisation of the simulated telemetry, whose role in this work is to justify that
the data on which the anomaly detection models are trained and evaluated behaves like a
powertrain.
